% ======================================================================
\section{Introduction}
% ======================================================================
Payment switches, core-banking platforms and the cloud services that host
them are prime targets. The financial sector had the second-highest
average breach cost of 17 industries in the IBM 2024 study, at
US\$6.08 million per incident~\cite{ibm2024breach}, and banks face DDoS,
ransomware, insider misuse and multi-stage
intrusions~\cite{asmar2024banks,islam2022ddosbank}. Signature-based
intrusion detection systems (IDS) miss unseen attacks, and a reactive
security operations centre (SOC) learns of an intrusion only after the damage
is done. \emph{Proactive} detection instead scores every flow as it is
observed and gives analysts a ranked, explained set of alerts.

Recent ML-based IDS use increasingly elaborate deep models: CNN--LSTM
hybrids~\cite{halbouni2022cnnlstm,altunay2023cnnlstm}, attention-augmented
recurrent networks~\cite{alashjaee2025attn} and
transformers~\cite{wu2022rtids,manocchio2024flowtransformer,akuthota2025iot}.
Deploying them in a regulated institution raises three issues that
benchmark accuracy does not capture.
\textit{(i) Tabular reality.} Flow records are tabular, where
gradient-boosted trees remain hard to beat~\cite{grinsztajn2022tree}, so a
deep detector must be compared with a tuned boosted-tree baseline. \textit{(ii) Privacy and compliance.} Flow
features such as transferred volumes, session durations and TTL or window
values act as \emph{quasi-identifiers} of customers and hosts. Under GDPR
Art.~25 (data protection by design)~\cite{gdpr2016} and PCI~DSS~\cite{pcidss4},
the detector and what it stores must not become a new leakage channel. \textit{(iii) Auditability.} Analysts need to know \emph{why} a flow
was flagged and how urgent it is~\cite{neupane2022xids}.

We present \textbf{HALX} (\underline{H}ybrid
\underline{A}ttention-\underline{L}STM + \underline{X}GBoost), which
addresses all three. Our contributions are:
\begin{itemize}
  \item A hybrid detector: a feature-token Attention-LSTM encoder learns a
        flow representation that an XGBoost head~\cite{chen2016xgboost}
        scores together with the privatized features, yielding a threat
        probability and three severity tiers.
  \item A privacy-preserving inference layer with \emph{formally stated}
        guarantees: Mondrian $k$-anonymity~\cite{lefevre2006mondrian} on
        quasi-identifiers and $(\varepsilon,\delta)$-DP preprocessing
        statistics with zCDP accounting~\cite{bun2016zcdp}.
  \item A reproducible evaluation on UNSW-NB15~\cite{moustafa2015unsw} and
        CICIDS2017~\cite{sharafaldin2018cicids} against seven baselines over
        three seeds, with ablation, privacy--utility, explanation and
        latency analyses. Code and results are released.\footnote{\url{https://github.com/mazharulmd/Threat-Detection-for-Financial-Systems}}
\end{itemize}

% ======================================================================
\section{Related Work}
\label{sec:related}
% ======================================================================
\textbf{Deep and hybrid IDS.} CNN--LSTM models combine convolutional
feature extraction with recurrent modelling and are a common baseline on
CICIDS2017 and UNSW-NB15~\cite{halbouni2022cnnlstm,altunay2023cnnlstm}.
Alashjaee~\cite{alashjaee2025attn} adds self-attention to a CNN--LSTM and
shows in an ablation that the attention layer helps.
Transformer-based detectors treat flow features as
tokens~\cite{wu2022rtids,long2024cloud,akuthota2025iot}. FlowTransformer
showed that the choice of input encoding and classification head matters more
than the transformer variant~\cite{manocchio2024flowtransformer}, and
BERT-style language models have also been adapted to
traffic~\cite{yang2025bert}. Closest to our work, Sajid et
al.~\cite{sajid2024hybrid} use XGBoost and a CNN as feature selectors in
front of an LSTM classifier. HALX reverses that order: the deep encoder
produces the representation and the boosted trees make the decision. This
follows evidence that trees are the stronger final learner on tabular
data~\cite{grinsztajn2022tree,gorishniy2021revisiting}.

\textbf{ML security in finance.} Surveys of digital banking identify DDoS,
ransomware, insider threats and phishing as the dominant threat
classes~\cite{asmar2024banks}. Islam et al.~\cite{islam2022ddosbank} detect
DDoS on a banking IoT monitoring system. Recent studies report SOC
deployments in core-banking systems~\cite{bhowmick2026cbs} and proactive
detection in emulated bank networks~\cite{mulea2026banking}. They stress
regulatory constraints, but none of them states a formal privacy guarantee
for the detector itself.

\textbf{Privacy-preserving IDS.} Most privacy work on IDS concerns
\emph{training}: federated learning~\cite{mcmahan2017fedavg} with
differentially private updates~\cite{ruzafa2023fl,gutti2025fl,vyas2024ppfl},
or DP-SGD~\cite{abadi2016dpsgd}. We address a complementary
problem that is often overlooked: what the deployed detector \emph{exposes},
namely generalized quasi-identifiers in its inputs and logs and the
normalization statistics shipped with the model. The two approaches can be
combined.

\textbf{Evaluation pitfalls.} CICIDS2017 contains labelling and flow-construction
errors~\cite{engelen2021troubleshooting}, and models trained on it transfer
poorly to related traffic~\cite{catillo2022transferability}. Arp et
al.~\cite{arp2022dos} catalogue pitfalls such as data snooping and
inappropriate baselines. We follow their recommendations: the test
partitions are never used for model selection, a tuned XGBoost baseline is
always included, and every number is reported as a mean over seeds.
Explainability is surveyed in~\cite{neupane2022xids}, and
\cite{barnard2022xai} pairs XGBoost with SHAP for NIDS. HALX instead
explains each alert through its intrinsic attention weights.
